\paragraph{The row at the chosen comparison.}
\label{par:linked-family-rows}
Fix a completed member \(q\) and its division pair from
Proposition~\ref{prop:completed-division-pair}. Choose a recursive
realization, a normalization of its real rows and the identification
between these realizations on the pair. Write \(P\) for the resulting point. The row,
periods, finite groups, coefficient field and numerical functions
at \(P\) are those of the pair. In particular its two level rows
are the ones in Section~\ref{subsec:completed-division}. The
following verifications hold for each compatible comparison on
the bundle. Thus they also hold if an additional condition, such
as the one on the evaluation grid, selects the comparison.

\begin{proposition}[The two rows at a compatible choice of comparison]\label{prop:linked-comparison-rows}
At every compatible comparison on the completed bundle, let \(q_0\) be projection to the sharp good row and \(P_{u,0}\) the reference pairing of word \(u\). The factorization is characterized by
\[
 F-F'\in\ker q_0\quad\Longrightarrow\quad P_{u,0}(\phi,F)=P_{u,0}(\phi,F')\quad\text{for all }\phi.
\] At either level the rational row and its integral
submodule have the corresponding divisor comparisons of
the original division pair, on all tuples used by the real
row. Any later choice of a comparison satisfying an extra
condition therefore still gives these rows and their pairings.
\end{proposition}

\begin{proof}
\proofstep{Use the sharp reference period to test the kernel of the good projection.}
Fix a word \(u\), and use its selected Poincar\'e operator from
Section~\ref{subsec:selected-completion-period}. Let \(R_P\) be the
rational row, \(R_P^\sharp\) its common sharp good row and
\(q_0:R_P\to R_P^\sharp\) the good projection. In the complex
realization write \(I_u\) for the column isometry and \(I_R\) for
the isometry on the good row at the refined level. Unfolding by
Lemma~\ref{lem:joint-period} gives, for the sharp tests of this
operator,
\[
 \iota\bigl(P_{u,0}(\phi,F)\bigr)
  =\mathcal P_{u,0}\bigl(I_u\phi, I_R(q_0 F)\bigr),
\]
where \(\mathcal P_{u,0}\) denotes the resulting complex theta
pairing.

The normalized factor is one at the reference tuple
of Section~\ref{sec:transport}. The reference tuple is obtained
from the full diagonal and the dimension of the ramified space
there, before applying either division estimate. If \(F-F'\in\ker q_0\),
the difference of the right sides vanishes for every \(\phi\).
The coefficient embedding is injective, so the vanishing difference implies
\[
 F-F'\in\ker q_0\quad\Longrightarrow\quad P_{u,0}(\phi,F)=P_{u,0}(\phi,F')\qquad(\text{every }\phi).
\] This defines the rational
reference pairing on the quotient without choosing an integral
representative.

\proofstep{Read the quotient pairing on the transported integral lattices.}

The rational pairing obtained this way is the one characterized by
the Gram equations in Proposition~\ref{prop:rational-period-factor}.
For each of the two levels, the rational row comparison identifies
the down-lift and the theta map in the same coordinates. To see what
is integral, let \(M_w\) be the localization of the glued lattice
at a place \(w\).

For the coordinate identification \(a\) coming
from the completed pair, its lattice and form are transported by
\[
 M_w\longmapsto aM_w,\qquad
 B_w(\phi,F)\longmapsto B_w(a^{-1}\phi,a^{-1}F).
\]
Here \(a\) denotes the respective coordinate map in each argument.
Thus the localized orders, lattice indices and dual indices are
computed on the identified modules. The reference form descends
because of the kernel calculation above. A nonreference tuple of
the real row occurs as the tuple of a place obtained by the inverse
place bijection of the host. Hence for each used tuple \(\mu\) one of two comparisons is available:
\[
 \mu=\mu_0\quad\text{or}\quad \mu=\mu(w)\ \text{for a place }w\text{ on the nonreference support}.
\]
The measured intersection comparison supplies the surjections
onto the good coordinates and the separation of the primitive types.
The equality of the bounded real realizations identifies the real
datum and its face in both calculations.
\label{cor:linked-family-rows-1}

On a nonreference summand, let \(c\) be an integral level coordinate
reading to an admitted source value and \(\pi_{\rm src}(c)\) its
projection to the lift source. Write \(x=\Theta\pi_{\rm src}(c)\)
in form coordinates and choose \(x=\sum_h a_h f_h\), with \(f_h\)
rational forms in the expansion model and \(a_h\) in the row field.
If \({\rm ev}_w\) denotes the local expansion map, the divisor theta
identity and the signed series comparison give the same value
\[
        {\rm ev}_w(x)=\sum_h\iota_w(a_h)\,{\rm ev}_w(f_h).
\]

To check independence of the expansion, subtract the two expressions for \(x\) in the extended space and use linearity:
\[
 \sum_h b_h f_h=0\quad\Longrightarrow\quad \sum_h b_h\,{\rm ev}_w(f_h)=0 .
\] The completed
division comparison transports \(\pi_{\rm src}\), the rational form
coordinates and \({\rm ev}_w\) through the update of periods and
sizes. Consequently the displayed equality can be checked either
before or after the update. In the chosen rational basis, evaluation on the row is the composite 
\[
 x\longmapsto(a_h)_h\longmapsto(\iota_w(a_h))_h\longmapsto\sum_h \iota_w(a_h)\,{\rm ev}_w(f_h).
\]
 Its restrictions to the integral expansion lattice and theta sublattice give the integral evaluation and the assigned series, respectively.
\label{cor:linked-family-rows-2}

Forget the grid in the arithmetic member of
Section~\ref{subsec:automorphic-realization}. The norm choice,
the signed series used for the exponent, the two type maps, the
base and local level groups, the kernel field and the selected
half primes remain those of this member. The changes in the
Petersson scalar, in the lift scalars and in the error have no
effect on the maps of finite ideles or integral host rings.

Fix a preliminary frame and its finite support data. If \(S\)
is an allowed support, put
\[
    M(S)=g_0g_1\!\!\prod_{\substack{q\text{ a residue cardinality}\\
                         q\text{ of a place in }S}}\!q ,
\]
with each distinct cardinality used once, as in the level of the
member. By the support bound the possible \(M(S)\) belong to a
finite list. If \(A_M\) generates the embedded ray field at level \(M\), one finite request suffices:
\[
 A=\bigcup_{M\text{ in the level list}} A_M,\qquad |A|\le\sum_M |A_M|<\infty .
\] The Artin comparison of
Section~\ref{subsec:automorphic-realization} identifies these ray
fields and the type and character values in the chosen embedding.
The two canonical frames on the same core give the same half-level
ideals as the divisor, including the half-error divisibility.

The integral field reading identifies the coefficient rings. After
passing to fractions it identifies the degree and absolute
discriminant used in the size bound with those of the recorded core.
The single-index support test uses the case with zero error and
height constants. These comparisons give the field, rank, series
and level requirements in the two runs. A coordinate isomorphism preserves each of the module equalities in the comparison. For submodules \(U,V\) in an identified row and its isomorphism \(b\),
\[
 U=V\quad\Longleftrightarrow\quad b(U)=b(V).
\]
\label{cor:linked-family-rows-3}

At infinity the intertwiner is defined on the prescribed lowest
compact type of the polynomial Weil representation. Regularity,
admissibility and finite multiplicity give the Harish--Chandra
modules. Apply Casselman--Wallach globalization to this intertwiner
and cuspidal rapid decay to its integrals, as in
Section~\ref{sec:transport}. Let \(T\) be the difference of two such intertwiners,
and let \(V_0\) be the span of the compact translates of the
distinguished lowest vectors. Equality on the test vectors gives
\(T(V_0)=0\).

At the Harish--Chandra module level,
\[
 V=U(\mathfrak g_{\mathbb C})V_0,\qquad
 T(Xv)=X T(v)=0
 \quad\bigl(X\in U(\mathfrak g_{\mathbb C}),\ v\in V_0\bigr).
\]
Thus \(T=0\) on the module and, by continuity, on its globalization. This explains the injective reading used for the geometric
action. The conjugate frame is read by the transported coefficient
map. When \(n_V\ge8fr+4\), the stable theta span and the form
isomorphism of Section~\ref{subsec:selected-cusp-reading}
apply to these weights. Since the normalized map and the extended
intertwiner agree on every vector of the common orthonormal basis,
they agree on the row. In particular evaluation and the two
linear readings assign the same form at either frame.

At a supported finite type, let \(e_{\rm good}\) be the central
idempotent of Section~\ref{subsec:automorphic-realization}. For a
vector \(f\) in its image and a finite type-group element \(g\),
\[
       e_{\rm good}(g f)=g(e_{\rm good}f)=g f .
\]
The chosen nonzero component therefore generates a module for
the supported split factors in the same constituent. Apply
Green's regular cuspidal character and the depth-zero type
correspondence to the factor actions. Normalized induction and
uniqueness of cuspidal support give the multiset of the local
parameter. The marked local Galois identification reads it at
the coefficient choice of \(P\). This applies to each required
label with the threshold on the family of members. The primitive
base lifting isomorphism is rational and bijective in its
phase-one coordinates and has the uniform positive rank \(D_0\).
The tuple map, level subgroups and base character therefore
refer to this same row on the two conjugate sides.
\label{cor:linked-family-rows-4}

\end{proof}

